\section{Fatto dopo la laurea}
\subsection{Jackknife fossili}
Jackknife dei fossili con i 186 geni BUSCO single copy e clock model 2, ovvero l'independent rates.

\begin{figure}[h!]
    \centering
    \includegraphics[width=1\textwidth]{Figures/Jackknife_scatterplot.png}
    \caption[Fossil Jackknife]{Dating of all nodes removing each time one node
}
    \label{fig:lsd2tree}
\end{figure}

\pagebreak
\subsection{gCF bassi}
gCF molto bassi a prescindere che siano nodi recenti o antichi:

\begin{figure}[h!]
    \centering
    \includegraphics[width=1\textwidth]{Figures/gCF_allspecies.tree.png}
    \caption[gCF Tree]{gCF tree with all species, each branch label is BS/gCF
}
    \label{fig:gCF_allspecies}
\end{figure}

\pagebreak
Essendo spesso Lynceus sp. a posizionarsi in topologie alternative, ho rifatto i gCF togliendo questa specie. La situazione non cambia:

\begin{figure}[h!]
    \centering
    \includegraphics[width=1\textwidth]{Figures/gCF_nolysp.tree.png}
    \caption[gCF Tree without Lynceys sp.]{gCF tree without Lynceus sp., each branch label is BS/gCF
}
    \label{fig:gCF_nolysp}
\end{figure}

\pagebreak
Questo si può vedere meglio dal seguente plot dei gCF coi due dataset.

\begin{figure}[h!]
    \centering
    \includegraphics[width=1\textwidth]{Figures/gCF.png}
    \caption[gCF Distribution Plot]{gCF distributions with the two datasets
}
    \label{fig:gCF_distributions}
\end{figure}

Non avendo ricevuto risposte certe riguardo il funzionamento dei gCF in caso di assenza di specie negli allineamenti dei singoli geni, ho provato a fare una filogenesi con gli allineamnti contenenti tutte le specie. Purtroppo sono solo 3, di conseguenza i gCF vengono per ogni nodo 0\%-33\%-66\%-100\%, non risolvendo il dubbio.

\pagebreak
\subsection{Topology test}
Visti i dubbi sui gCF e sulle topologie discordi, ho fatto due gene tree per ogni gene, uno senza constraint e uno con constraint la topologia generale. Ho poi fatto dei topology test a coppie gene per gene per vedere quali geni rifiutassero la topologia constraint.
\\Risultato:
\begin{itemize}
    \item 28 geni non rifiutano la topologia (detti per comodità "concordant genes")
    \item 159 geni rifiutano la topologia (detti "discordant genes")
\end{itemize}
Ho quindi fatto due allineamenti concatenati e rifatto il dating. Come si vede dal grafico seguente, non si evidenziano particolari differenze con il dating con tutti i geni insieme.

\begin{figure}[h!]
    \centering
    \includegraphics[width=1\textwidth]{Figures/All_Conc_Disc_scatterplot.png}
    \caption[AllConcDisc Scatterplot]{Comparison of the date estimates for each node with the three datasets: all genes, concordant genes, discordant genes.
}
    \label{fig:topologytest_allspecies}
\end{figure}
Visto quanto osservato con Lynceus, ho deciso di rifare il processo togliendolo dagli allineamenti. Questo ha aumentato notevolmente il numero di concordant genes.
\begin{itemize}
    \item 181 concordant genes
    \item 6 discordant genes
\end{itemize}
Purtroppo però ancora non ho ancora potuto vedere se/come cambia il dating con questo subset.